Ein Propellerwagen (Masse \mWO) steht auf einer langen, ebenen Schiene; die Reibung ist vernachlässigbar. Zur Zeit $t = 0$ wird der Propeller eingeschaltet und schiebt den Wagen mit einer konstanten Kraft an. Zur Zeit $t_1 = \tEO$ hat der Wagen die Geschwindigkeit $v_1 = \vEO$, und der Propeller wird abgeschaltet.

\begin{center}
\begin{tikzpicture}[>=Stealth]
 \draw[thick] (-0.5,0) -- (6.5,0);
 \foreach \x in {-0.4,-0.1,...,6.4} \draw (\x,0) -- ++(-0.15,-0.15);
 \draw[fill=gray!25] (2,0.3) rectangle (3.6,0.9);
 \draw[fill=white] (2.35,0.17) circle (0.15) (3.25,0.17) circle (0.15);
 \draw[thick] (2.25,0.9) -- (2.25,1.35);
 \draw[fill=gray!50] (2.25,1.35) ellipse (0.08 and 0.4);
 \draw[->, very thick] (3.6,0.6) -- (4.9,0.6) node[right] {$\vec F$};
\end{tikzpicture}
\end{center}

\begin{abcliste}
\abc
\DuSie{Skizziere}{Skizzieren Sie} den Graphen der Geschwindigkeit gegen die Zeit von $t = 0$ bis $t_2 = \tZO$.

\abc
Wie schnell fährt der Wagen zur Zeit $t_2$, und wie weit ist er dann vom Start entfernt?

\abc
Wie gross ist die Kraft des Propellers?

\abc
Lea meint: „Nach dem Abschalten wird der Wagen langsam langsamer, weil sein Schwung aufgebraucht wird.“ Was \DuSie{meinst du}{meinen Sie} dazu?
\end{abcliste}
